% =============================================================================
% Anexo D · Reproducibilidad: datos congelados (D4), paridad Python–R (D6) y orden de ejecución
% Fuentes: decisiones.qmd (D4, D6), CLAUDE.md (comandos), data/results/final/R_paridad_python.csv.
% Las huellas md5 son identificadores de D4 (clase «Identificadores» de la lista de excepciones).
% =============================================================================
\chapter{Reproducibilidad}
\label{ap:reproducibilidad}

Este anexo documenta cómo se reproduce el análisis. Complementa la
\cref{sec:met-implementacion}, que describe la implementación y la paridad.

\section{Datos congelados}
\label{sec:ap-datos}

Por la decisión D4, la tesis usa la versión de los datos de \fechaInicioNiveles{} a
\fechaFinDatos{} versionada en el repositorio el 16 de octubre de 2026. El código comprueba la
huella md5 de cada archivo antes de leerlo y se detiene si no coincide
(\cref{tab:md5}). Cualquier actualización de los datos requiere una nueva decisión, porque el BCRP
revisa el ajuste estacional de algunas series y sus valores cambian con la fecha de descarga.

\begin{table}[htbp]
\centering
\small
\caption{Archivos de datos congelados y su huella md5 (decisión D4)}
\label{tab:md5}
\begin{tabular}{@{}>{\raggedright\arraybackslash}p{0.5\textwidth}l@{}}
\toprule
Archivo (en \path{data/processed/}) & md5 \\
\midrule
\path{dataset_niveles_mensual.csv} & \texttt{8088c28cf277419388d9562722faa1ad} \\
\path{dataset_transformado_mensual.csv} & \texttt{8c49f3e55608f2651be3c1dee4236194} \\
\path{dataset_niveles_estandarizado_mensual.csv} & \texttt{153638578212ecc820ec9223087c080a} \\
\path{dataset_estandarizado_mensual.csv} & \texttt{43a01dead7ffb96e8b3b7e4d2799f15d} \\
\path{series_bcrp_consolidado.csv} & \texttt{61e6a6bc92eee23c6c476f096d402edb} \\
\bottomrule
\end{tabular}

\smallskip
\begin{minipage}{0.95\textwidth}\footnotesize
Fuente: \texttt{decisiones.qmd}, decisión D4. Las exclusiones de series (crédito total del sistema
financiero e IPC subyacente por D1; liquidez M2 por D16) se aplican al leer los archivos, que no
se modifican.
\end{minipage}
\end{table}

\section{Paridad entre Python y R}
\label{sec:ap-paridad}

Por la decisión D6, todo el análisis se programa en Python y en R, y la versión en R es la
referencia. La tabla de paridad del análisis principal
(\path{data/results/final/R_paridad_python.csv}) compara cada resultado en las dos lenguas con
la tolerancia numérica declarada en cada fila; en la versión vigente, coinciden
\filasParidadCoinciden{}
de \filasParidad{} comparaciones. El anexo de
robustez tiene su propia tabla (\path{data/results/final/anexo/R_paridad_python.csv}). Los
sorteos aleatorios los hace R con semilla fija y los escribe en archivos que Python lee, de modo
que la paridad no depende de los generadores de números aleatorios.

\section{Orden de ejecución}
\label{sec:ap-orden}

Desde la raíz del repositorio, la cadena vigente se reproduce en este orden:

\begin{enumerate}[leftmargin=*]
  \item notebooks de Python \path{05_pca}, \path{05c_anexo_no_diferenciar},
    \path{06_epca}, \path{06c_aede_grilla}, \path{07b_nulo_estructura},
    \path{07c_hfc_clases}, \path{09a_datacion_pca} y \path{09b_walkforward_pca}, con
    \texttt{jupyter nbconvert -{}-execute};
  \item el documento final en R, \path{R/ciclo_financiero_final.Rmd}, con
    \texttt{rmarkdown::render}, que escribe los archivos \path{R_*.csv} y la tabla de paridad;
  \item \texttt{python3} \path{herramientas/generar_cifras.py}, que regenera
    \path{tesis_latex/cifras.tex}; la opción \texttt{-{}-comprobar} verifica que esté al día;
  \item \texttt{latexmk -lualatex -cd} \path{tesis_latex/tesis.tex}, que compila este documento.
\end{enumerate}

Las pruebas automáticas del repositorio (\texttt{pytest} \path{tests/}) verifican las funciones de
\path{src/}, incluida la paridad del port de \texttt{aede} y \texttt{clahfac}.

\section{Estados anteriores recuperables}
\label{sec:ap-estados}

Los estados del análisis anteriores a las decisiones vigentes se conservan en el control de
versiones: el tag \texttt{v0.1-transformadas} (series transformadas como base, anterior a D1), el
tag \texttt{v0.2-niveles} (anterior a D2) y la rama \path{archivo/v0.3-trimestral} (frecuencia
trimestral, anterior a D3). El estado con la liquidez M2 y en dos vías, anterior a D15 y D16, está en el
commit \texttt{2be3665}; el tag \texttt{v0.4-17series-dos-vias} que debe marcarlo aún no se ha
creado (\cref{sec:pend-operativos}).
